% Author: Cole Rodney Migisha (first draft written 16 Nov from the repository; check and extend).
\section{Visual and Audio Pipelines}
\label{sec:pipelines}

\subsection{Mouth crops}
The visual input is a sequence of grey $96 \times 96$ mouth crops at 25~frames/s
(\texttt{lipscribe/\allowbreak{}preprocessing/\allowbreak{}crop.py}). MediaPipe Face Mesh finds the lip landmarks in every
frame; a square box around them, enlarged by a 25~\% margin, is smoothed over time with a
centred moving average, and frames where the face is lost take the box of their neighbours.
The crop is padded by repeating edge pixels when the box leaves the frame, and a per-frame
validity mask records where Face Mesh succeeded. The model therefore always receives 75 frames
per GRID clip, and knows which of them show a real face.

During training the crops are augmented with a random shift of up to 4 pixels, a horizontal
flip with probability 0.5, a scale change of up to 10~\%, a random erased patch (probability
0.5, up to 15~\% of the area) and two time masks of up to 6 frames each. Evaluation and the
application use the crops unchanged.

\subsection{Audio features}
Audio is resampled to 16~kHz mono and converted on the GPU to 80-band log-Mel features with a
25~ms window, a 10~ms hop and a 512-point FFT, then standardised with the mean and standard
deviation of the training set (\texttt{lipscribe/\allowbreak{}preprocessing/\allowbreak{}features.py}). Experiment X5
compared these features with 39 MFCCs and with log-Mel plus MFCC on the audio-only baseline;
log-Mel plus MFCC lowered the mean validation WER from 7.72~\% to 6.90~\%
(\texttt{reports/\allowbreak{}w3\_\allowbreak{}08\_\allowbreak{}audio\_\allowbreak{}features\_\allowbreak{}decision.md}).
P4 nevertheless keeps the 80 log-Mel features (\texttt{configs/\allowbreak{}p\_\allowbreak{}init\_\allowbreak{}b1\_\allowbreak{}b2\_\allowbreak{}v2.yaml}). P4 was
trained on 5~Nov from the B1 audio encoder of 3~Nov, before the X5 runs of 10~Nov, and its
checkpoint was frozen on 16~Nov. Switching features would have meant retraining B1 and then P4, and
the X5 gain was measured on validation for the audio-only model only. Every GRID model compared in
Section~\ref{sec:results} therefore uses log-Mel 80, so the comparisons between them are unaffected.

\subsection{Encoders}
The audio encoder applies two strided 1-D convolutions ($80 \rightarrow 256 \rightarrow 256$,
kernel~3, stride~2), which reduce the 100~frames/s features to the video rate of 25~frames/s,
followed by a two-layer bidirectional LSTM with 256 units per direction. Both visual encoders
map $[B, 1, T, 96, 96]$ crops to $[B, T, 512]$. V1 stacks three 3-D convolution blocks (32, 64
and 96 channels) before a linear layer and a two-layer BiLSTM. V2 follows the lip-reading
front end of Stafylakis and Tzimiropoulos~\cite{stafylakis2017}: one $5 \times 7 \times 7$ 3-D
convolution, a ResNet-18 trunk applied to each frame, and the same BiLSTM. V2 learned faster
and generalised better, so model P4 uses it.

\subsection{Real-world pipeline}
\label{sec:pipelines-real}
The GRID encoders do not carry over to open-vocabulary speech. Model~P's audio encoder learned
51~words, and a fused model built from the same parts and trained on about 28~hours of film and
TED clips (Section~\ref{sec:data}) still made 72~\% character errors from clean audio, no fewer
than its audio branch alone (\texttt{reports/\allowbreak{}flix\_\allowbreak{}v4/\allowbreak{}decision\_\allowbreak{}log.md}). We therefore rebuilt the
real-world model the way Model~P was built: each stream is first trained on its own from a
pretrained starting point, and fusion starts from the finished streams. Two rules applied
throughout. Pretrained weights serve only as starting points, so every reported model is our own
fine-tuning. No Whisper-based audio encoder is used, because the training labels come from
Whisper~\cite{radford2023whisper}, and such a model would be scored against its own output.

\textit{Audio stream.} WavLM Base+~\cite{chen2022wavlm} is fine-tuned with a CTC loss on the film
and TED clips, and its 50~frames/s output is pooled to 25~frames/s, so that both streams share
one time axis. The first run overfitted: its validation loss rose from epoch~3 while the training
loss kept falling. The run we kept was therefore regularised more strongly: time masking 0.15,
channel masking 0.05, LayerDrop 0.1, head dropout 0.2, half the learning rate and weight decay
0.05. This lowered the mean validation WER from 30.1~\% to 29.2~\%.

\textit{Visual stream.} The visual encoder is the lip-reading front end and Conformer of
Auto-AVSR~\cite{ma2023}. Fine-tuned on our first crops, it fell behind the unmodified model,
because those tight lip boxes differ from the crops Auto-AVSR was trained on. We re-implemented
its crop with MediaPipe landmarks (\texttt{lipscribe/\allowbreak{}preprocessing/\allowbreak{}avsr\_\allowbreak{}crop.py}). The landmarks
are smoothed over 12~frames, eight stable points are aligned to Auto-AVSR's mean face at
$256 \times 256$, and a $96 \times 96$ grey patch is cut around the mouth. On 142~TED validation
clips, the unmodified model's WER fell from 67.9~\% with our old crops to 45.3~\% with the new
ones, so every clip was cropped again. The encoder also keeps Auto-AVSR's pretrained word-piece
CTC layer, rather than a new character layer, so all our real-world models predict the same
5,000 word pieces.

\textit{Fusion.} We tried three fusion designs in turn, and kept the last.
\begin{enumerate}
\item \textit{Cross-attention fusion, as in Model~P, with a new character head.} A probe showed
that it ignored the lips: giving each clip another clip's video left WER unchanged, and the gate
leaned about 90~\% to audio. The pretrained audio dominated from the first epoch.
\item \textit{Lip-anchored fusion} ($f = v + \mathrm{MLP}([a; v])$, where MLP is a multilayer perceptron, zero-initialised, word-piece
heads). It started as the visual model and learned to add the audio. It used the lips: shuffling
the video raised its TED WER at $-10$~dB from 36.2~\% to 93.7~\%. However, it also over-trusted
them, so it was 9.5~points worse than audio alone on clean TED speech, and fine-tuning could not
remove the built-in lip path.
\item \textit{Reliability-gated fusion} (\texttt{V4FusionGated} in
\texttt{lipscribe/\allowbreak{}models/\allowbreak{}pretrained.py}). Both streams are frozen. Each frame's features are
projected and combined as $f = g \odot a' + (1 - g) \odot v' + \mathrm{MLP}([a'; v'])$, where the
gate $g$ has one value per channel (768~channels). The gate sees both features and each stream's
own confidence, which is the entropy and the blank probability of that stream's CTC output.
Transformer layers and a new word-piece head follow.
\end{enumerate}
The final model, \textit{ws\_v5}, is a warm start of the third design. Its audio side is the whole
frozen path of an \textit{audio-only twin}: the same architecture trained with the lips always
hidden. The gate starts 95~\% open to audio, the new layer starts as the identity, and the head is
copied from the twin. At the first step the warm start therefore equals the twin (TED validation
WER 20.9~\% against 20.8~\%), and training can only add what the lips contribute. The twin,
\textit{twin\_v5}, is also the baseline for every real-world comparison, since it differs from
ws\_v5 only by the lips.

\textit{Clip windows (v5).} Our first real-world clips were cut tightly around their words. On the
TED test set the twin got a clip's first or last word wrong 22~\% of the time, against 11~\% for
the words inside it, and widening only the test audio by 0.32~s per side lowered the edge-word
error to 19~\%. Every clip was therefore cut again with 8~frames (0.32~s) more on each side.
Its mouth crops and audio come from the same source window, and its label is every Whisper
large-v3-turbo word with at least half its duration inside that window, read from a cut with
1.5~s of context on each side. The four stages (lips, twin, lip-anchored fusion, warm start) were
then retrained on these clips, with their recipes unchanged.

\textit{Decoding.} The CTC output is decoded by beam search with a word 4-gram language model
(LM) built with KenLM~\cite{heafield2011}. One LM uses our training transcripts only. The other adds 40~million
words of LibriSpeech's LM text~\cite{panayotov2015}, with our transcripts weighted five times. No
validation or test talk appears in either. The LM, its weight $\alpha$, the word bonus $\beta$ and
the beam width were chosen per model on the film and TED validation sets at every noise level.
The LibriSpeech LM won for every model. For ws\_v5 the choice is $\alpha = 0.2$, $\beta = 1.0$
and width~32 (\texttt{reports/\allowbreak{}flix\_\allowbreak{}v5/\allowbreak{}final\_\allowbreak{}decoder.md}).
